\documentclass[11pt,a4paper]{article}

% ============================================
% PACCHETTI
% ============================================
\usepackage[utf8]{inputenc}
\usepackage[T1]{fontenc}
\usepackage[italian]{babel}
\usepackage{geometry}
\usepackage{hyperref}
\usepackage{xcolor}
\usepackage{listings}
\usepackage{tikz}
\usepackage{booktabs}
\usepackage{longtable}
\usepackage{array}
\usepackage{fancyhdr}
\usepackage{titlesec}
\usepackage{float}
\usepackage{enumitem}

% ============================================
% CONFIGURAZIONE
% ============================================
\geometry{margin=2.5cm}
\hypersetup{
    colorlinks=true,
    linkcolor=blue,
    urlcolor=cyan,
    pdftitle={Exam Management Component - Technical Documentation}
}

% Colori
\definecolor{codegreen}{rgb}{0,0.6,0}
\definecolor{codegray}{rgb}{0.5,0.5,0.5}
\definecolor{codepurple}{rgb}{0.58,0,0.82}
\definecolor{backcolour}{rgb}{0.95,0.95,0.92}
\definecolor{darkblue}{rgb}{0.0,0.0,0.6}
\definecolor{deepred}{rgb}{0.6,0,0}
\definecolor{deepgreen}{rgb}{0,0.5,0}

% Stile codice
\lstdefinestyle{mystyle}{
    backgroundcolor=\color{backcolour},
    commentstyle=\color{codegreen},
    keywordstyle=\color{magenta},
    numberstyle=\tiny\color{codegray},
    stringstyle=\color{codepurple},
    basicstyle=\ttfamily\footnotesize,
    breaklines=true,
    numbers=left,
    numbersep=5pt
}
\lstset{style=mystyle}

% Tikz
\usetikzlibrary{shapes.geometric, arrows, positioning}

% Stile titoli
\titleformat{\section}{\Large\bfseries\color{darkblue}}{\thesection}{1em}{}
\titleformat{\subsection}{\large\bfseries\color{darkblue}}{\thesubsection}{1em}{}

% Header
\pagestyle{fancy}
\fancyhf{}
\fancyhead[L]{\leftmark}
\fancyhead[R]{Exam Management Component}
\fancyfoot[C]{\thepage}
\renewcommand{\headrulewidth}{0.4pt}

% Titolo
\title{
    \Huge\textbf{Exam Management Component}\\[0.5em]
    \Large\textit{Documentazione Tecnica Completa}\\[1em]
    \large Silvi - Exam Planning \& Study Organization
}
\author{
    \textbf{Extra Tracker Development Team}\\
    \texttt{Versione 1.0 - Febbraio 2026}
}
\date{\today}

\begin{document}

\maketitle
\thispagestyle{empty}

\vfill

\begin{abstract}
\noindent
Documentazione tecnica del componente \textbf{Exam Management} di Silvi. Il sistema permette agli studenti di organizzare i propri esami, associare mazzi di flashcard, tracciare il progresso di studio e gestire le scadenze. Include funzionalita' di pianificazione, filtri per stato (attivo/urgente/completato), statistiche di padronanza e integrazione con il sistema di studio.
\end{abstract}

\vfill
\newpage

\tableofcontents
\newpage

% ============================================
% SEZIONE 1: INTRODUZIONE
% ============================================
\section{Introduzione}\label{sec:intro}

\subsection{Panoramica}

Il componente \textbf{Exam Management} gestisce il ciclo di vita completo degli esami:

\begin{itemize}[leftmargin=*]
    \item \textbf{Creazione Esami}: Titolo, descrizione, data scadenza
    \item \textbf{Associazione Decks}: Collegamento mazzi di studio all'esame
    \item \textbf{Tracciamento Progresso}: Percentuale padronanza, carte da ripassare
    \item \textbf{Gestione Scadenze}: Badge urgenti, notifiche scadenza imminente
    \item \textbf{Stati Esame}: Active, Urgent, Completed, Passed, Failed, Archived
    \item \textbf{Statistiche}: Deck count, total cards, mastery percent, card distribution
\end{itemize}

\subsection{Relazioni con Altri Componenti}

\begin{center}
\begin{tikzpicture}[
    box/.style={rectangle, rounded corners, draw=darkblue, fill=blue!5, 
                minimum width=2.8cm, minimum height=1cm, align=center, font=\small},
    arrow/.style={->, >=stealth, thick}
]
    \node[box] (exam) at (0,0) {\textbf{Exam}\\Management};
    \node[box] (deck) at (5,0) {\textbf{Deck}\\Flashcards};
    \node[box] (study) at (10,0) {\textbf{Study}\\Session};
    \node[box] (auth) at (0,-2.5) {\textbf{Auth}\\User};
    \node[box] (folder) at (5,-2.5) {\textbf{Folder}\\Organization};
    
    \draw[arrow] (exam) -- (deck) node[midway,above,font=\tiny] {examId};
    \draw[arrow] (deck) -- (study) node[midway,above,font=\tiny] {cards};
    \draw[arrow] (auth) -- (exam) node[midway,left,font=\tiny] {user};
    \draw[arrow] (folder) -- (deck) node[midway,below,font=\tiny] {folderId};
    \draw[arrow, dashed] (exam) -- (study) node[midway,right,font=\tiny] { aggregated stats};
\end{tikzpicture}
\end{center}

\subsection{Struttura File}

\begin{lstlisting}[language=bash]
# BACKEND
server/
|-- models/
|   |-- Exam.js                 # Schema esame con outcome
|-- controllers/
|   |-- examsController.js      # CRUD esami
|-- routes/
|   |-- exams.js                # /api/exams/*

# FRONTEND
src/features/study/
|-- types/
|   |-- exam.ts                 # TypeScript definitions
|-- services/
|   |-- examService.ts          # API client
|-- components/
|   |-- Exams/
|   |   |-- ExamsView.tsx       # Lista esami principale
|   |   |-- ExamCard.tsx        # Card singolo esame
|   |   |-- ExamsFilters.tsx    # Filtri e ordinamento
|   |   |-- utils/
|   |   |   |-- examIcons.tsx   # Icone dinamiche per materia
\end{lstlisting}

% ============================================
% SEZIONE 2: DATA MODEL
% ============================================
\section{Data Model}\label{sec:model}

\subsection{Exam Schema}

\begin{lstlisting}[language=JavaScript, caption={Exam Model - server/models/Exam.js}]
const outcomeSchema = new mongoose.Schema({
    grade: {
        type: Number,
        default: null,
        min: [0, 'Il voto non puo essere negativo'],
    },
    date: {
        type: Date,
        default: null,
    },
    notes: {
        type: String,
        default: '',
        trim: true,
        maxlength: [1000, 'Le note non possono superare 1000 caratteri'],
    },
    difficulties: [{  // Tag difficolta' segnalate
        type: String,
        trim: true,
    }],
}, { _id: false });

const examSchema = new mongoose.Schema({
    title: { 
        type: String, 
        required: [true, 'Il titolo e obbligatorio'],
        trim: true,
        maxlength: [120, 'Il titolo non puo superare 120 caratteri'],
    },
    description: {
        type: String,
        default: '',
        trim: true,
        maxlength: [1000, 'La descrizione non puo superare 1000 caratteri'],
    },
    deadline: {
        type: Date,
        required: [true, 'La data dell\'esame e obbligatoria'],
    },
    status: {
        type: String,
        enum: ['active', 'passed', 'failed', 'archived', 'completed'],
        default: 'active',
    },
    outcome: {
        type: outcomeSchema,
        default: null,
    },
}, { timestamps: true });

// Indici
examSchema.index({ user: 1, deadline: 1 });

// Multi-tenancy
examSchema.plugin(multiTenancyPlugin, {
    tenantField: 'user',
    required: true,
});
\end{lstlisting}

\subsection{TypeScript Types}

\begin{lstlisting}[language=TypeScript, caption={Exam Types - src/features/study/types/exam.ts}]
export type ExamStatus = 'active' | 'passed' | 'failed' | 'archived' | 'completed';

export interface ExamOutcome {
    grade?: number;
    date: string;
    notes?: string;
    difficulties?: string[];  // Per recovery plan
}

export interface Exam {
    id: string;
    title: string;
    description?: string;
    deadline: string;  // ISO date
    status: ExamStatus;
    outcome?: ExamOutcome | null;
    createdAt?: string;
    updatedAt?: string;
}

export interface CreateExamPayload {
    title: string;
    description?: string;
    deadline: string;
}

export interface UpdateExamPayload {
    title?: string;
    description?: string;
    deadline?: string;
    status?: ExamStatus;
    outcome?: ExamOutcome | null;
}
\end{lstlisting}

% ============================================
% SEZIONE 3: BACKEND API
% ============================================
\section{Backend API}\label{sec:backend}

\subsection{Controller}

\begin{lstlisting}[language=JavaScript, caption={Exams Controller - server/controllers/examsController.js}]
const Exam = require('../models/Exam');
const Deck = require('../models/Deck');

// GET /api/exams - Tutti gli esami
const getAllExams = asyncHandler(async (req, res) => {
    const userId = req.tenantScope.userId;
    const exams = await Exam.find({ user: userId })
        .sort({ createdAt: -1 });
    res.json({ success: true, data: exams });
});

// GET /api/exams/:id - Singolo esame
const getExamById = asyncHandler(async (req, res) => {
    const { id } = req.params;
    const userId = req.tenantScope.userId;
    
    const exam = await Exam.findOne({ _id: id, user: userId });
    if (!exam) throw AppError.notFound('Esame non trovato');
    
    res.json({ success: true, data: exam });
});

// POST /api/exams - Crea esame
const createExam = asyncHandler(async (req, res) => {
    const { title, description, deadline, status, outcome } = req.body;
    const userId = req.tenantScope.userId;
    
    if (!title?.trim()) throw AppError.validation('Il titolo e obbligatorio');
    if (!deadline) throw AppError.validation('La data e obbligatoria');
    
    const exam = await Exam.create({
        title: title.trim(),
        description: description?.trim() || '',
        deadline: new Date(deadline),
        status: status || 'active',
        outcome: outcome || null,
        user: userId,
    });
    
    res.status(201).json({ success: true, data: exam });
});

// PATCH /api/exams/:id - Aggiorna esame
const updateExam = asyncHandler(async (req, res) => {
    const { id } = req.params;
    const updates = req.body;
    const userId = req.tenantScope.userId;
    
    const exam = await Exam.findOne({ _id: id, user: userId });
    if (!exam) throw AppError.notFound('Esame non trovato');
    
    // Aggiorna campi
    if (updates.title !== undefined) exam.title = updates.title.trim();
    if (updates.description !== undefined) exam.description = updates.description?.trim() || '';
    if (updates.deadline !== undefined) exam.deadline = new Date(updates.deadline);
    if (updates.status !== undefined) exam.status = updates.status;
    if (updates.outcome !== undefined) exam.outcome = updates.outcome;
    
    await exam.save();
    res.json({ success: true, data: exam });
});

// DELETE /api/exams/:id - Elimina esame
const deleteExam = asyncHandler(async (req, res) => {
    const { id } = req.params;
    const userId = req.tenantScope.userId;
    
    const exam = await Exam.findOne({ _id: id, user: userId });
    if (!exam) throw AppError.notFound('Esame non trovato');
    
    // IMPORTANTE: Rimuovi examId dai deck associati
    await Deck.updateMany(
        { user: userId, examId: id },
        { $unset: { examId: 1 } }
    );
    
    await exam.deleteOne();
    res.json({ success: true, message: 'Esame eliminato' });
});
\end{lstlisting}

\subsection{Routes}

\begin{lstlisting}[language=JavaScript, caption={Exam Routes - server/routes/exams.js}]
const express = require('express');
const router = express.Router();
const { requireAuth } = require('../middleware/auth');
const { tenantContext } = require('../middleware/tenantContext');
const examsController = require('../controllers/examsController');

// Proteggi tutte le routes
router.use(requireAuth);
router.use(tenantContext({ required: true }));

router.get('/', examsController.getAllExams);
router.post('/', examsController.createExam);
router.get('/:id', examsController.getExamById);
router.patch('/:id', examsController.updateExam);
router.delete('/:id', examsController.deleteExam);

module.exports = router;
\end{lstlisting}

% ============================================
% SEZIONE 4: FRONTEND
% ============================================
\section{Frontend Architecture}\label{sec:frontend}

\subsection{Service Layer}

\begin{lstlisting}[language=TypeScript, caption={Exam Service - src/features/study/services/examService.ts}]
import { apiClient, type ApiResponse } from '../../../shared/services/apiClient';
import type { Exam, CreateExamPayload, UpdateExamPayload } from '../types/exam';

const unwrap = <T>(response: ApiResponse<T>, fallbackMessage: string): T => {
    if (!response.success || response.data === undefined) {
        throw new Error(response.error?.message || response.message || fallbackMessage);
    }
    return response.data;
};

const examService = {
    async getAll(): Promise<Exam[]> {
        const response = await apiClient.get<Exam[]>('/exams');
        return unwrap(response, 'Errore nel recupero degli esami');
    },

    async getById(id: string): Promise<Exam> {
        const response = await apiClient.get<Exam>(`/exams/${id}`);
        return unwrap(response, `Errore nel recupero esame ${id}`);
    },

    async create(payload: CreateExamPayload): Promise<Exam> {
        const response = await apiClient.post<Exam>('/exams', payload);
        return unwrap(response, 'Errore nella creazione dell\'esame');
    },

    async update(id: string, payload: UpdateExamPayload): Promise<Exam> {
        const response = await apiClient.patch<Exam>(`/exams/${id}`, payload);
        return unwrap(response, `Errore nell'aggiornamento esame ${id}`);
    },

    async delete(id: string): Promise<void> {
        const response = await apiClient.delete<null>(`/exams/${id}`);
        unwrap(response, `Errore nell'eliminazione esame ${id}`);
    },
};

export default examService;
\end{lstlisting}

\subsection{ExamCard Component}

\begin{lstlisting}[language=TypeScript, caption={ExamCard - src/features/study/components/Exams/ExamCard.tsx}]
interface ExamCardProps {
    exam: Exam;
    deckCount: number;
    totalCards: number;
    dueCards: number;
    masteryPercent: number;
    decks?: Deck[];
    onClick: () => void;
    onDelete?: (examId: string) => void;
    onReactivate?: (examId: string) => void;
}

export const ExamCard: React.FC<ExamCardProps> = ({
    exam, deckCount, totalCards, dueCards, masteryPercent,
    decks = [], onClick, onDelete, onReactivate
}) => {
    // Calcola giorni alla scadenza
    const deadlineDate = new Date(exam.deadline);
    const daysUntil = Math.ceil(
        (deadlineDate.getTime() - Date.now()) / (1000 * 60 * 60 * 24)
    );
    const isUrgent = daysUntil <= 7 && daysUntil >= 0;
    const isOverdue = daysUntil < 0;
    
    // Distribuzione carte per status
    const distribution = useMemo(() => {
        const examDecks = decks.filter(d => d.examId === exam.id);
        const allCards = examDecks.flatMap(d => d.cards || []);
        
        return {
            new: allCards.filter(c => c.status === 'new').length,
            learning: allCards.filter(c => c.status === 'learning').length,
            review: allCards.filter(c => c.status === 'review').length,
            mastered: allCards.filter(c => c.status === 'mastered').length,
        };
    }, [decks, exam.id]);
    
    // Configurazione status badge
    const statusConfig = useMemo(() => {
        switch (exam.status) {
            case 'passed':
                return { label: 'Superato', icon: CheckCircle, color: 'text-emerald-400' };
            case 'failed':
                return { label: 'Non Superato', icon: XCircle, color: 'text-rose-400' };
            case 'archived':
                return { label: 'Archiviato', icon: Archive, color: 'text-slate-400' };
            default:
                if (isOverdue) return { label: 'Scaduto', icon: AlertCircle, color: 'text-red-400' };
                if (isUrgent) return { label: 'Urgente', icon: AlertCircle, color: 'text-orange-400' };
                return { label: 'In Corso', icon: Sparkles, color: 'text-primary-400' };
        }
    }, [exam.status, isOverdue, isUrgent]);
    
    return (
        <motion.div
            whileHover={{ scale: 1.01, y: -2 }}
            onClick={onClick}
            className={`
                rounded-2xl border overflow-hidden cursor-pointer
                ${isUrgent || isOverdue
                    ? 'border-orange-500/40 bg-gradient-to-br from-orange-500/10'
                    : 'border-primary-500/30 bg-gradient-to-br from-primary-500/10'
                }
            `}
        >
            {/* Header */}
            <div className="p-5">
                <div className="flex items-start gap-4 mb-4">
                    {/* Icona dinamica */}
                    <div className={`p-3 rounded-xl ${statusConfig.bgColor}`}>
                        <ExamIcon className={`w-6 h-6 ${statusConfig.color}`} />
                    </div>
                    
                    <div className="flex-1">
                        <h3 className="font-bold text-white text-lg">{exam.title}</h3>
                        <span className={`
                            px-2 py-0.5 rounded-full text-[10px] font-bold
                            ${statusConfig.bgColor} ${statusConfig.color}
                        `}>
                            <statusConfig.icon className="w-3 h-3" />
                            {statusConfig.label}
                        </span>
                    </div>
                </div>
                
                {/* Barra distribuzione carte */}
                {totalCards > 0 && (
                    <div className="mb-4">
                        <div className="h-2 bg-white/10 rounded-full overflow-hidden flex">
                            {distribution.new > 0 && (
                                <div className="h-full bg-blue-500" 
                                     style={{ width: `${(distribution.new/totalCards)*100}%` }} />
                            )}
                            {distribution.learning > 0 && (
                                <div className="h-full bg-amber-500" 
                                     style={{ width: `${(distribution.learning/totalCards)*100}%` }} />
                            )}
                            {distribution.review > 0 && (
                                <div className="h-full bg-orange-500" 
                                     style={{ width: `${(distribution.review/totalCards)*100}%` }} />
                            )}
                            {distribution.mastered > 0 && (
                                <div className="h-full bg-emerald-500" 
                                     style={{ width: `${(distribution.mastered/totalCards)*100}%` }} />
                            )}
                        </div>
                    </div>
                )}
                
                {/* Stats Grid */}
                <div className="grid grid-cols-2 gap-3">
                    <div className="p-2.5 rounded-xl bg-white/5 border border-white/10">
                        <Layers className="w-4 h-4 text-primary-400" />
                        <p className="text-lg font-bold text-white">{deckCount}</p>
                        <p className="text-[10px] text-white/50">Mazzi</p>
                    </div>
                    <div className="p-2.5 rounded-xl bg-white/5 border border-white/10">
                        <Target className="w-4 h-4 text-emerald-400" />
                        <p className="text-lg font-bold text-white">{masteryPercent}%</p>
                        <p className="text-[10px] text-white/50">Padronanza</p>
                    </div>
                </div>
            </div>
            
            {/* Footer - Scadenza */}
            <div className="px-5 py-3 border-t border-white/10 bg-white/[0.03]">
                <div className="flex items-center justify-between">
                    <div className="flex items-center gap-2.5">
                        <Calendar className={`w-3.5 h-3.5 ${isOverdue ? 'text-red-400' : 'text-white/60'}`} />
                        <p className={`text-sm font-semibold ${isOverdue ? 'text-red-400' : 'text-white/80'}`}>
                            {isOverdue 
                                ? `Scaduto (${Math.abs(daysUntil)}gg fa)`
                                : daysUntil === 0 ? 'Oggi!'
                                : daysUntil === 1 ? 'Domani'
                                : `${daysUntil} giorni`
                            }
                        </p>
                    </div>
                    <ArrowRight className="w-4 h-4 text-primary-400" />
                </div>
            </div>
        </motion.div>
    );
};
\end{lstlisting}

\subsection{ExamsView - Lista Principale}

\begin{lstlisting}[language=TypeScript, caption={ExamsView - Lista e Filtri}]
export const ExamsView: React.FC<ExamsViewProps> = ({
    decks, onCreateExam, onExamClick, onDeckUpdate, ...
}) => {
    const [exams, setExams] = useState<Exam[]>([]);
    const [isLoading, setIsLoading] = useState(true);
    const [searchQuery, setSearchQuery] = useState('');
    const [sortBy, setSortBy] = useState<ExamSortOption>('recent');
    const [filter, setFilter] = useState<ExamFilterOption>('all');
    
    // Carica esami
    useEffect(() => {
        loadExams();
    }, []);
    
    const loadExams = async () => {
        try {
            setIsLoading(true);
            const allExams = await examService.getAll();
            setExams(allExams);
        } catch (err) {
            setError(err.message);
        } finally {
            setIsLoading(false);
        }
    };
    
    // Calcola statistiche per esame
    const getExamStats = (examId: string) => {
        const examDecks = decks.filter(d => d.examId === examId);
        const totalCards = examDecks.reduce((sum, d) => sum + (d.totalCards || 0), 0);
        const dueCards = examDecks.reduce((sum, d) => sum + (d.dueCount || 0), 0);
        const masteredCards = examDecks.reduce(
            (sum, d) => sum + (d.cards?.filter(c => c.status === 'mastered').length || 0), 
            0
        );
        const masteryPercent = totalCards > 0 
            ? Math.round((masteredCards / totalCards) * 100) 
            : 0;
        
        return { deckCount: examDecks.length, totalCards, dueCards, masteryPercent };
    };
    
    // Filtra e ordina
    const filteredAndSortedExams = useMemo(() => {
        let filtered = [...exams];
        const now = Date.now();
        
        // Filtro per ricerca
        if (searchQuery.trim()) {
            const query = searchQuery.toLowerCase();
            filtered = filtered.filter(exam =>
                exam.title.toLowerCase().includes(query) ||
                exam.description?.toLowerCase().includes(query)
            );
        }
        
        // Filtro per stato
        if (filter === 'urgent') {
            filtered = filtered.filter(exam => {
                if (exam.status !== 'active') return false;
                const deadline = new Date(exam.deadline).getTime();
                const daysUntil = Math.ceil((deadline - now) / (1000 * 60 * 60 * 24));
                return daysUntil <= 7 && daysUntil >= 0;
            });
        } else if (filter === 'upcoming') {
            filtered = filtered.filter(exam => {
                if (exam.status !== 'active') return false;
                const deadline = new Date(exam.deadline).getTime();
                const daysUntil = Math.ceil((deadline - now) / (1000 * 60 * 60 * 24));
                return daysUntil > 7;
            });
        } else if (filter === 'completed') {
            filtered = filtered.filter(exam => 
                ['completed', 'passed', 'failed', 'archived'].includes(exam.status)
            );
        } else if (filter === 'all') {
            filtered = filtered.filter(exam => exam.status === 'active');
        }
        
        // Ordinamento
        filtered.sort((a, b) => {
            switch (sortBy) {
                case 'recent':
                    return new Date(b.createdAt).getTime() - new Date(a.createdAt).getTime();
                case 'deadline':
                    return new Date(a.deadline).getTime() - new Date(b.deadline).getTime();
                case 'name':
                    return a.title.localeCompare(b.title);
                case 'mastery': {
                    const aStats = getExamStats(a.id);
                    const bStats = getExamStats(b.id);
                    return bStats.masteryPercent - aStats.masteryPercent;
                }
                case 'cards': {
                    const aStats = getExamStats(a.id);
                    const bStats = getExamStats(b.id);
                    return bStats.totalCards - aStats.totalCards;
                }
                default: return 0;
            }
        });
        
        return filtered;
    }, [exams, searchQuery, filter, sortBy, decks]);
    
    // Render
    return (
        <div className="space-y-6">
            <ExamsFilters 
                searchQuery={searchQuery}
                onSearchChange={setSearchQuery}
                sortBy={sortBy}
                onSortChange={setSortBy}
                filter={filter}
                onFilterChange={setFilter}
            />
            
            <div className="grid grid-cols-1 md:grid-cols-2 lg:grid-cols-3 gap-4">
                {filteredAndSortedExams.map(exam => {
                    const stats = getExamStats(exam.id);
                    return (
                        <ExamCard
                            key={exam.id}
                            exam={exam}
                            {...stats}
                            onClick={() => onExamClick(exam.id)}
                        />
                    );
                })}
            </div>
        </div>
    );
};
\end{lstlisting}

% ============================================
% SEZIONE 5: BUSINESS LOGIC
% ============================================
\section{Business Logic}\label{sec:business}

\subsection{Calcolo Statistiche Aggregated}

Le statistiche di un esame sono calcolate aggregando i dati di tutti i deck associati:

\begin{lstlisting}[language=TypeScript]
// Calcolo statistiche esame
const getExamStats = (examId: string, decks: Deck[]) => {
    // Filtra deck associati all'esame
    const examDecks = decks.filter(d => d.examId === examId);
    
    // Conteggio mazzi
    const deckCount = examDecks.length;
    
    // Conteggio totale carte
    const totalCards = examDecks.reduce(
        (sum, deck) => sum + (deck.totalCards || deck.cards?.length || 0), 
        0
    );
    
    // Carte da ripassare (in scadenza)
    const dueCards = examDecks.reduce(
        (sum, deck) => sum + (deck.dueCount || 0), 
        0
    );
    
    // Carte padroneggiate
    const masteredCards = examDecks.reduce((sum, deck) => {
        return sum + (deck.cards?.filter(c => c.status === 'mastered').length || 0);
    }, 0);
    
    // Percentuale padronanza
    const masteryPercent = totalCards > 0 
        ? Math.round((masteredCards / totalCards) * 100) 
        : 0;
    
    return { deckCount, totalCards, dueCards, masteryPercent };
};
\end{lstlisting}

\subsection{Filtraggio per Scadenza}

\begin{lstlisting}[language=TypeScript]
// Logica filtri scadenza
const now = Date.now();

// Urgente: entro 7 giorni
const isUrgent = (exam: Exam) => {
    if (exam.status !== 'active') return false;
    const deadline = new Date(exam.deadline).getTime();
    const daysUntil = Math.ceil((deadline - now) / (1000 * 60 * 60 * 24));
    return daysUntil <= 7 && daysUntil >= 0;
};

// Prossimi: dopo 7 giorni
const isUpcoming = (exam: Exam) => {
    if (exam.status !== 'active') return false;
    const deadline = new Date(exam.deadline).getTime();
    const daysUntil = Math.ceil((deadline - now) / (1000 * 60 * 60 * 24));
    return daysUntil > 7;
};

// Scaduto: deadline passata
const isOverdue = (exam: Exam) => {
    const deadline = new Date(exam.deadline).getTime();
    return deadline < now;
};
\end{lstlisting}

\subsection{Visualizzazione Scadenza}

\begin{lstlisting}[language=TypeScript]
// Formattazione human-readable scadenza
const formatDeadline = (deadline: string): string => {
    const deadlineDate = new Date(deadline);
    const daysUntil = Math.ceil(
        (deadlineDate.getTime() - Date.now()) / (1000 * 60 * 60 * 24)
    );
    
    if (daysUntil < 0) {
        return `Scaduto (${Math.abs(daysUntil)}gg fa)`;
    }
    if (daysUntil === 0) return 'Oggi!';
    if (daysUntil === 1) return 'Domani';
    if (daysUntil <= 7) return `${daysUntil} giorni (urgente)`;
    return `${daysUntil} giorni`;
};
\end{lstlisting}

% ============================================
% SEZIONE 6: API ENDPOINTS
% ============================================
\section{API Endpoints}\label{sec:endpoints}

\begin{longtable}{|p{4cm}|p{1.5cm}|p{8cm}|}
\hline
\textbf{Endpoint} & \textbf{Method} & \textbf{Descrizione} \\
\hline
\endfirsthead
\texttt{/api/exams} & GET & Lista tutti gli esami dell'utente \\
\hline
\texttt{/api/exams/:id} & GET & Dettaglio singolo esame \\
\hline
\texttt{/api/exams} & POST & Crea nuovo esame \\
\hline
\texttt{/api/exams/:id} & PATCH & Aggiorna esame (titolo, deadline, status, outcome) \\
\hline
\texttt{/api/exams/:id} & DELETE & Elimina esame e rimuove associazioni deck \\
\hline
\end{longtable}

% ============================================
% APPENDICE
% ============================================
\appendix
\section{Sequence Diagram}

\textbf{Creazione Esame}:
\begin{enumerate}
    \item User clicca "Crea Esame" nel frontend
    \item Modale apre con form (titolo, descrizione, deadline)
    \item User submit → \texttt{examService.create(payload)}
    \item POST /api/exams con JWT e CSRF token
    \item Controller valida input
    \item Exam.create() con userId dal tenantScope
    \item Ritorna 201 Created
    \item Frontend aggiorna lista esami
\end{enumerate}

\textbf{Eliminazione Esame}:
\begin{enumerate}
    \item User clicca menu → Elimina su ExamCard
    \item Conferma modale mostra warning: "X deck saranno disassociati"
    \item User conferma → \texttt{examService.delete(id)}
    \item DELETE /api/exams/:id
    \item Controller: Deck.updateMany() rimuove examId dai deck associati
    \item Controller: exam.deleteOne()
    \item Ritorna 200 Success
    \item Frontend rimuove esame dalla lista (optimistic UI)
\end{enumerate}

\section{Glossario}

\begin{description}
    \item[Mastery Percent] Percentuale carte padroneggiate (status = mastered)
    \item[Due Cards] Carte in scadenza per ripasso (nextReviewDate <= now)
    \item[Exam Outcome] Risultato esame (grade, date, notes, difficulties)
    \item[Urgent] Scadenza entro 7 giorni
    \item[Recovery Plan] Piano di ripasso basato sulle difficulties segnalate
\end{description}

\end{document}
